
%-------------------------
% Author : Isin Altinkaya
% Website: https://github.com/isinaltinkaya/cv
%------------------------

\documentclass[letterpaper,10.5pt]{article}
\usepackage{latexsym}
\usepackage[empty]{fullpage}
\usepackage{titlesec}
\usepackage{marvosym}
\usepackage[usenames,dvipsnames]{color}
\usepackage{verbatim}
\usepackage{enumitem}
\usepackage{fancyhdr}
\usepackage{graphicx}
\usepackage{bold-extra}
\usepackage[utf8]{inputenc}  % input encoding
\usepackage{mathtools} % align text

%icons
\usepackage{orcidlink}


\newcommand\sbullet[1][.5]{\mathbin{\vcenter{\hbox{\scalebox{#1}{$\bullet$}}}}}

\usepackage{pdfpages}


\usepackage{xcolor}
\definecolor{isored}{HTML}{980000}


% \newcommand\ytl[2]{
% \parbox[b]{8em}{\hfill{\color{cyan}\bfseries\sffamily #1}~$\cdots\cdots$~}\makebox[0pt][c]{$\bullet$}\vrule\quad \parbox[c]{4.5cm}{\vspace{7pt}\color{red!40!black!80}\raggedright\sffamily #2.\\[7pt]}\\[-3pt]}
% \begin{table}
% \ytl{1947}{AT and T Bell Labs develop the idea of cellular phones}
% \ytl{1968}{Xerox Palo Alto Research Centre envisage the `Dynabook'}
% \ytl{1971}{Busicom 'Handy-LE' Calculator}
% \ytl{1973}{First mobile handset invented by Martin Cooper}
% \ytl{1978}{Parker Bros. Merlin Computer Toy}
% \ytl{1981}{Osborne 1 Portable Computer}
% \ytl{1982}{Grid Compass 1100 Clamshell Laptop}
% \ytl{1983}{TRS-80 Model 100 Portable PC}
% \ytl{1984}{Psion Organiser Handheld Computer}
% \ytl{1991}{Psion Series 3 Minicomputer}
% \end{table}


\usepackage[pdftex]{hyperref}

\hypersetup{
    colorlinks = true,
    linkcolor = blue,
    anchorcolor = blue,
    citecolor = blue,
    filecolor = blue,
    urlcolor = blue,
}

\usepackage{tikz}
\newcommand{\pie}[1]{%
    \begin{tikzpicture}[scale=0.15]
        \draw (0,0) circle (1);
        \fill[black] (0,0) -- (0:1) arc (0:#1:1) -- cycle;
    \end{tikzpicture}
}

% Adjust margins
\addtolength{\oddsidemargin}{-0.375in}
\addtolength{\evensidemargin}{-0.375in}
\addtolength{\textwidth}{1in}
\addtolength{\topmargin}{-.5in}
\addtolength{\textheight}{1.0in}
\urlstyle{same}
\raggedbottom
\raggedright
\setlength{\tabcolsep}{0in}

% Sections formatting
\titleformat{\section}{
  \vspace{-4pt}\bfseries\scshape\raggedright\large\color{isored}
}{}{1em}{}[\color{isored}{\titlerule} \vspace{-3pt}]

% Custom commands
\newcommand{\cvItem}[2]{
  \item\small{
    \textbf{#1}{: #2 \vspace{-2pt}}
  }
}

\newcommand{\cvSubheading}[5]{
  \vspace{-1pt}\item
    \begin{tabular*}{0.97\textwidth}{l@{\extracolsep{\fill}}r}
      \textbf{#1} & #2 \\
      {\small#3} & {\small #4} \\
    \end{tabular*}\vspace{3pt} \\
    #5
}

\newcommand{\cvSubItem}[2]{\cvItem{#1}{#2}\vspace{-4pt}}

%\renewcommand{\labelitemii}{$\circ$}
%\renewcommand{\labelitemii}{$\rightarrow$}


\newcommand{\cvSubHeadingListStart}{\begin{itemize}[leftmargin=*]}
\newcommand{\cvSubHeadingListEnd}{\end{itemize}}
\newcommand{\cvItemListStart}{\begin{itemize}}
\newcommand{\cvItemListEnd}{\end{itemize}\vspace{-5pt}}
%
%\newcommand\textline[4][t]{%
%  \par\smallskip\noindent\parbox[#1]{.333\textwidth}{\raggedright\texttt#2}%
%  \parbox[#1]{.333\textwidth}{\centering#3}%
%  \parbox[#1]{.333\textwidth}{\raggedleft\texttt{#4}}\par\smallskip%
%}

%-------------------------------------------
%%%%%%  CV STARTS HERE  %%%%%%%%%%%%%%%%%%%%

\begin{document}

%----------HEADING-----------------

\href{http://isinaltinkaya.com/}{isinaltinkaya.com} 
\color{isored}
\hfill \begin{Huge} Isin Altinkaya \end{Huge} \\
\color{black}

\href {https://github.com/isinaltinkaya}{$\vcenter{\hbox{\includegraphics[scale=0.04]{icons/github_PNG23.png}/isinaltinkaya}}$} \hfill Oester Voldgade 5-7, 1350 \\
\href {https://x.com/IsinAltinkaya}{$\vcenter{\hbox{\includegraphics[scale=0.04]{icons/Twitter_and_X_logos.png}}}$/isinaltinkaya}
\hfill Copenhagen K, Denmark \\
\href{mailto:isinaltinkaya@gmail.com}{isinaltinkaya@gmail.com} $\sbullet[.5]$ \href{mailto:isin.altinkaya@sund.ku.dk}{isin.altinkaya@sund.ku.dk} \hfill \href{https://scholar.google.com/citations?user=DHpwS6oAAAAJ&hl=en&oi=ao}{$\vcenter{\hbox{\includegraphics[scale=0.08]{icons/google-scholar_orig_full.png}}}$} $\sbullet[.5]$
\href{https://orcid.org/0000-0002-6364-3332}{$\vcenter{\hbox{\includegraphics[scale=0.08]{icons/orcid-orig_full.png}}}$} $\sbullet[.5]$ 
\href{https://www.linkedin.com/in/isinaltinkaya/}{$\vcenter{\hbox{\includegraphics[scale=0.08]{icons/linkedin-orig-full.png}}}$}


% \href {https://github.com/isinaltinkaya}{GitHub/isinaltinkaya} 
% \hfill Oester Voldgade 5-7, 1350 \\
% \href {https://x.com/IsinAltinkaya}{X(Twitter)/isinaltinkaya}
% \hfill Copenhagen K, Denmark \\
% \href{mailto:isinaltinkaya@gmail.com}{isinaltinkaya@gmail.com} $\sbullet[.5]$ \href{mailto:isin.altinkaya@sund.ku.dk}{isin.altinkaya@sund.ku.dk} \hfill \href{https://scholar.google.com/citations?user=DHpwS6oAAAAJ&hl=en&oi=ao}{\includegraphics[scale=0.08]{icons/google-scholar_orig_full.png}} $\sbullet[.5]$
% \href{https://orcid.org/0000-0002-6364-3332}{{\includegraphics[scale=0.08]{icons/orcid-orig_full.png}} $\sbullet[.5]$ 
% \href{https://www.linkedin.com/in/altinkayaisin/}{\includegraphics[scale=0.07]{icons/linkedin-orig-full.png}}}



%I\c{s}{\i}n Alt{\i}nkaya

\section{Nachweis des Grades der Behinderung}

\begin{flushleft}

\medskip
Sehr geehrtes Einstellungskomitee,\\
\vskip 0.2in


im Anschluss an Ihre Bitte, den Nachweis meines Grades der Behinderung hochzuladen, übersende ich Ihnen hiermit die erforderlichen Unterlagen. Ich wurde in Dänemark, einem EU-Mitgliedstaat, diagnostiziert und gelte gemäß der dänischen Gesetzgebung als behindert. Als Nachweis füge ich mein Ledsagekort bei – die dänische Behinderten-Identifikationskarte, ausgestellt von den Danske Handicaporganisationer (Behindertenorganisation Dänemark).

\medskip
re
Bitte beachten Sie, dass in Dänemark Behinderungen nicht prozentual bewertet werden, sondern ausschließlich in die Kategorien "behindert" und "nicht behindert" eingeteilt werden. Daher liegt mir kein spezifischer Prozentsatz meines Behinderungsgrades vor. Allerdings ordnet \href{https://www.gesetze-im-internet.de/versmedv/BJNR241200008.html}{die Versorgungsmedizin-Verordnung (VersMedV, BJNR 241200008)} meine Diagnose in den Bereich GdS 50--70 ein.

\medskip

Sollten Sie eine ärztliche Bescheinigung benötigen, kann ich Ihnen auf Anfrage eine Note meines behandelnden Arztes in Dänemark zukommen lassen. Da ich derzeit nicht in Deutschland ansässig bin und keiner deutschen Krankenkasse angehöre, kann ich momentan leider kein deutsches Dokument vorlegen. Sobald ich jedoch nach Deutschland umziehe, werde ich Ihnen gerne ein entsprechendes deutsches Attest nachreichen.

\vskip 0.2in
Mit freundlichen Grüßen,\\
Isin Altinkaya

%\noindent\makebox[\linewidth]{\rule{0.9\paperwidth}{0.1pt}}
\section{Proof of the degree of disability}

\medskip
Dear Hiring Committee,\\
\vskip 0.2in

In response to your request to upload proof of my degree of disability, I hereby submit the required documents. I was diagnosed in Denmark, an EU member state, and am recognized as disabled in accordance with Danish law. As evidence, I am attaching my Ledsagekort – the Danish disability identification card issued by Danske Handicaporganisationer (Disabled People's Organisation Denmark).

\medskip

Please note that in Denmark, disabilities are not evaluated in percentages but are classified solely as either "disabled" or "not disabled." Therefore, I do not have a specific percentage rating of my disability. However, \href{https://www.gesetze-im-internet.de/versmedv/BJNR241200008.html}{the Versorgungsmedizin-Verordnung (VersMedV, BJNR 241200008)} classifies my diagnosis within the range of GdS 50--70.

\medskip

Should you require a doctor's note, I can provide a statement from my treating physician in Denmark upon request. As I am currently not residing in Germany and am not affiliated with a German health insurance provider, I am unable to furnish a German document at this time. However, I will gladly provide one once I relocate to Germany.

\vskip 0.2in
Sincerely,\\
Isin Altinkaya

\includepdf[pages=-,pagecommand={},width=\textwidth]{IsinAltinkaya_Ledsagekort.pdf}

\end{flushleft}

\end{document}  